%% report_at2_phase_field.tex
%% Standalone technical report: AT2 phase-field fracture of growing biofilms
%% Compile: lualatex (x3) report_at2_phase_field.tex
%%
\documentclass[11pt, a4paper]{article}

% ── Packages ─────────────────────────────────────────────────────────────────
\usepackage[T1]{fontenc}
\usepackage[utf8]{inputenc}
\usepackage{lmodern}
\usepackage{microtype}
\usepackage[margin=25mm]{geometry}
\usepackage{amsmath, amssymb, bm}
\usepackage{graphicx}
\usepackage{booktabs}
\usepackage{caption}
\usepackage{subcaption}
\usepackage[colorlinks=true, linkcolor=blue!60!black, citecolor=blue!60!black,
            urlcolor=blue!60!black]{hyperref}
\usepackage{cleveref}

% ── Macros ───────────────────────────────────────────────────────────────────
\newcommand{\Gc}{G_c}
\newcommand{\ac}{\alpha_c}
\newcommand{\psiE}{\psi_e}
\newcommand{\divop}{\operatorname{div}}
\newcommand{\Hfield}{H}
\newcommand{\um}{$\mu$m}        % micro-meter without siunitx
\newcommand{\invsec}{s$^{-1}$}

% ── Title ────────────────────────────────────────────────────────────────────
\title{\textbf{AT2 Phase-Field Fracture of Growing Biofilms}\\[4pt]
       \large Fully Coupled 2-D Q4 FEM Implementation and Parameter Sensitivity}
\author{K.\ Nishioka \\ Keio University / IKM Stuttgart}
\date{\today}

% ─────────────────────────────────────────────────────────────────────────────
\begin{document}
\maketitle
\tableofcontents
\bigskip

% ─────────────────────────────────────────────────────────────────────────────
\section{Introduction}
\label{sec:intro}

Biofilm detachment from surfaces is a clinically relevant fracture process:
commensal and dysbiotic oral biofilms differ in growth kinetics
($k_{\mathrm{eff}}^{b}$, s$^{-1}$), and the pathogen
\textit{Porphyromonas gingivalis} (\textit{Pg}) colonises preferentially at the
biofilm--substrate interface.
This report documents the implementation of a variational AT2 phase-field fracture
model for a biaxially constrained thin biofilm growing on a rigid substrate,
together with a material-parameter sensitivity study.

% ─────────────────────────────────────────────────────────────────────────────
\section{Variational Formulation}
\label{sec:theory}

\subsection{Energy functional (AT2, Miehe et al.\ 2010)}

The total free energy of the coupled system \((d, \bm{u})\) reads
\begin{equation}
  \mathcal{E}(d, \bm{u})
  = \int_\Omega \Bigl[
      g(d)\,\psiE
      + \frac{\Gc}{2\ell}\bigl(d^2 + \ell^2\,|\nabla d|^2\bigr)
    \Bigr]\,\mathrm{d}\Omega,
  \label{eq:energy}
\end{equation}
where $g(d) = (1-d)^2 + k_{\mathrm{res}}$ is the quadratic degradation function
($k_{\mathrm{res}} = 10^{-8}$ prevents singularity at $d=1$),
$\Gc$ [J/m$^2$] is the fracture energy, and $\ell$ [m] is the
phase-field length scale controlling the diffuse crack width.

\subsection{Plane-strain mechanical energy}

The film is modelled in plane strain ($\varepsilon_{yy} = 0$ in the
out-of-plane direction).  With isotropic growth eigenstrain
$\bm{\varepsilon}^0 = \alpha(x,z,t)\,\bm{I}$ the mechanical strain is
$\bm{\varepsilon}^m = \bm{\varepsilon} - \bm{\varepsilon}^0$ and
\begin{equation}
  \psiE = \tfrac12\,\bm{\sigma}^0 : \bm{\varepsilon}^m,
  \qquad
  \bm{\sigma}^0 = \mathbb{C} : \bm{\varepsilon}^m,
  \label{eq:psi}
\end{equation}
where $\mathbb{C}$ is the plane-strain elasticity tensor
($\lambda = \nu E / [(1+\nu)(1-2\nu)]$, $\mu = E/[2(1+\nu)]$).

\subsection{Damage evolution PDE}

The Euler--Lagrange equation for $d$ with the irreversibility constraint
$\dot{d} \ge 0$ (via the history field $\Hfield = \max_{s\le t}\psiE$) is
\begin{equation}
  -\Gc\ell\,\Delta d
  + \Bigl(\frac{\Gc}{\ell} + 2\Hfield\Bigr)d = 2\Hfield,
  \qquad
  \nabla d \cdot \bm{n} = 0 \text{ on } \partial\Omega.
  \label{eq:pde}
\end{equation}

\subsection{Critical eigenstrain}

Setting $\psiE = \Gc/(2\ell)$ in the prescribed-eigenstrain limit
(infinite film, no mechanical coupling) gives the analytic critical strain
\begin{equation}
  \ac
  = \sqrt{\frac{\Gc(1-\nu)}{2E\ell}}.
  \label{eq:alphac}
\end{equation}
This is the key scaling parameter for the sensitivity study in
\cref{sec:sensitivity}.

% ─────────────────────────────────────────────────────────────────────────────
\section{Numerical Method}
\label{sec:method}

\subsection{Staggered scheme}

At each time step, given $d^{k-1}$:
\begin{enumerate}
  \item Solve 2-D plane-strain elasticity with degraded stiffness $g(d^{k-1})$:
        \begin{equation}
          \divop\bigl[g(d)\,\mathbb{C}:(\bm{\varepsilon}-\bm{\varepsilon}^0)\bigr]
          = \bm{0}, \quad \bm{u}=\bm{0}\ \text{at substrate},\
          \bm{t}=\bm{0}\ \text{elsewhere}.
        \end{equation}
  \item Compute $\psiE$ from the actual (undegraded) strain field;
        update $\Hfield = \max(\Hfield, \psiE)$.
  \item Solve \cref{eq:pde} for $d^k$; enforce irreversibility
        $d^k \ge d^{k-1}$.
\end{enumerate}

\subsection{Spatial discretisation}

\begin{itemize}
  \item \textbf{Mechanics}: bilinear Q4 FEM on $(N_x-1)\times(N_z-1)$ cells,
        one-point Gauss quadrature (reduced integration avoids volumetric
        locking at $\nu=0.45$).
  \item \textbf{Phase field}: 5-point FD on $N_x\times N_z$ nodes;
        Neumann ghost-node operator (Kronecker product).
  \item \textbf{Mesh adequacy}: $h \le \ell/2$ required; at
        $N_x=120$, $N_z=80$ and $\ell=5$\,\um: $h/\ell = 0.17$.
  \item \textbf{Linear algebra}: \texttt{scipy.sparse.linalg.spsolve} with CSR
        matrices; Dirichlet BCs via in-place row zeroing (no format conversion).
\end{itemize}

% ─────────────────────────────────────────────────────────────────────────────
\section{Simulation Cases and Eigenstrain Profiles}
\label{sec:cases}

Three cases model distinct biofilm phenotypes
(\cref{tab:cases}).

\begin{table}[h]
  \centering
  \caption{Simulation cases.  $k_{\mathrm{eff}}^{\mathrm{CH}} = \ac\cdot1.6/120$\,s$^{-1}$,
           $k_{\mathrm{eff}}^{\mathrm{DH}} = k_{\mathrm{eff}}^{\mathrm{CH}}/k_{\mathrm{ratio}}$.}
  \label{tab:cases}
  \begin{tabular}{lll}
    \toprule
    Label & $k_{\mathrm{eff}}$ & Profile $k(x,z)$ \\
    \midrule
    CH & $k_{\mathrm{eff}}^{\mathrm{CH}}$ & z-gradient (higher at film surface) \\
    DH-base & $k_{\mathrm{eff}}^{\mathrm{DH}}$ & z-gradient \\
    DH$+$\textit{Pg} & $k_{\mathrm{eff}}^{\mathrm{DH}}$ & Pg substrate (concentrated at $z=0$) \\
    \bottomrule
  \end{tabular}
\end{table}

\noindent
The theoretical growth-rate ratio
$k_{\mathrm{ratio}} = [k_{\mathrm{eff}}^b(\text{CH})/k_{\mathrm{eff}}^b(\text{DH})]
= 6.44^{1/2.68} \approx 2.004$
directly predicts $t_{\mathrm{crit}}(\text{DH-base})/t_{\mathrm{crit}}(\text{CH})$.

% ─────────────────────────────────────────────────────────────────────────────
\section{Coupled Simulation Results}
\label{sec:coupled}

\cref{fig:coupled} shows the fully-coupled AT2 results at
$N_x=120$, $N_z=80$ ($h/\ell=0.17$, domain $100\,\um\times60\,\um$).

\begin{figure}[htb]
  \centering
  \includegraphics[width=\linewidth]{phase_field_at2_2d_thesis.pdf}
  \caption{Fully-coupled AT2 phase-field fracture (\texttt{inferno} colormap:
           dark~=~intact, bright~=~fractured; white dashed: $d=0.5$ fracture front).
           \textbf{Top row} (CH): damage initiates at substrate corners and
           propagates as a substrate-level delamination band.
           \textbf{Bottom row} (DH$+$\textit{Pg}): substrate-concentrated eigenstrain
           drives uniform delamination across the full width.
           Effect~1: $t_{\mathrm{crit}}^{\mathrm{DH}}/t_{\mathrm{crit}}^{\mathrm{CH}}
           = 2.004$ (exact match with $k_{\mathrm{ratio}}$).
           Effect~2 varies with domain width (\cref{sec:wide}).}
  \label{fig:coupled}
\end{figure}

\subsection{Key numerical results}

\begin{table}[h]
  \centering
  \caption{Fracture time $t_{\mathrm{crit}}$ ($d_{\max}>0.5$) and final maximum
           damage for the baseline domain ($W=100$\,\um).}
  \label{tab:results}
  \begin{tabular}{lccc}
    \toprule
    Case & $t_{\mathrm{crit}}$ [s] & $d_{\max}^{\mathrm{final}}$ & Ratio \\
    \midrule
    CH       & 56.0 & 0.955 & -- \\
    DH-base  & 112.2 & 0.955 & $2.004$ \\
    DH$+$\textit{Pg} & 38.1 & 0.958 & $0.680$ \\
    \bottomrule
  \end{tabular}
\end{table}

Effect~2 ($t_{\mathrm{crit}}^{\mathrm{DH}+Pg}/t_{\mathrm{crit}}^{\mathrm{CH}} = 0.680$):
despite a $2\times$ slower growth rate, the \textit{Pg}-substrate eigenstrain
concentration drives delamination \emph{faster} than CH in the narrow domain
where edge effects amplify substrate stresses.

\subsection{Wide domain ($W = 300$\,\um)}
\label{sec:wide}

\begin{table}[h]
  \centering
  \caption{Same cases with $W=300$\,\um, $N_x=300$, $N_z=80$.}
  \label{tab:wide}
  \begin{tabular}{lccc}
    \toprule
    Case & $t_{\mathrm{crit}}$ [s] & $d_{\max}^{\mathrm{final}}$ & Ratio \\
    \midrule
    CH       & 27.0 & 0.966 & -- \\
    DH-base  & 54.1 & 0.966 & $2.004$ \\
    DH$+$\textit{Pg} & 34.1 & 0.956 & $1.262$ \\
    \bottomrule
  \end{tabular}
\end{table}

In the wider domain, the interior is shielded from edge effects.
Effect~2 reverses ($1.262 > 1$): the slower DH growth rate now
dominates over the \textit{Pg} substrate concentration.
\emph{Domain width is therefore a key model parameter.}

% ─────────────────────────────────────────────────────────────────────────────
\section{Parameter Sensitivity ($\ell$, $G_c$)}
\label{sec:sensitivity}

\cref{fig:params} shows how the critical eigenstrain $\ac$ and the fracture
time $t_{\mathrm{crit}}$ depend on the length scale $\ell$ and fracture energy $\Gc$
at a fixed physical growth rate $\dot\alpha = K_{\mathrm{eff}}^{\mathrm{CH,base}}
= 3.13\times10^{-4}$\,\invsec\ (prescribed-eigenstrain mode, CH case,
$N_x=120$, $N_z=80$).

\begin{figure}[htb]
  \centering
  \includegraphics[width=\linewidth]{param_study_thesis.pdf}
  \caption{\textbf{Left}: $\ac = \sqrt{\Gc(1-\nu)/(2E\ell)}$ over the
           $(\Gc, \ell)$ grid.
           \textbf{Centre}: resulting fracture time $t_{\mathrm{crit}}$
           (numbers in cells, seconds); $t_{\mathrm{crit}} \propto \ac$.
           \textbf{Right}: $d_{\max}(t)$ curves for varying $\ell$ at fixed
           $\Gc = 10^{-5}$\,J/m$^2$.
           Dashed box marks the baseline parameters.}
  \label{fig:params}
\end{figure}

\subsection{Scaling law}

From \cref{eq:alphac} and \cref{eq:pde} and the prescribed-eigenstrain limit:
\begin{equation}
  t_{\mathrm{crit}} \propto \ac = \sqrt{\frac{\Gc(1-\nu)}{2E\ell}},
  \label{eq:scaling}
\end{equation}
confirmed numerically over the full $4\times4$ grid (\cref{fig:params}, centre).

\subsection{Length-scale effect}

Larger $\ell$ lowers $\ac$ and therefore \emph{accelerates} fracture at fixed
growth rate --- a counterintuitive result that arises because a wider diffuse
damage zone reaches the AT2 activation threshold at a smaller eigenstrain.
In the limit $\ell \to 0$ (sharp crack), $\ac \to \infty$ and the model
predicts no fracture for finite eigenstrains, consistent with the sharp-crack
Griffith theory.

% ─────────────────────────────────────────────────────────────────────────────
\section{Summary and Open Points}
\label{sec:summary}

\begin{itemize}
  \item Fully-coupled AT2/Q4-FEM staggered solver implemented and verified:
        $t_{\mathrm{crit}}^{\mathrm{DH}}/t_{\mathrm{crit}}^{\mathrm{CH}}
        = 2.004$ (exact, \cref{eq:alphac} and \cref{eq:scaling}).
  \item Dominant fracture mode is \emph{substrate delamination}
        (compressive lateral eigenstrain drives interfacial failure), not
        channeling.
  \item Effect~2 (\textit{Pg} substrate) is geometry-dependent: $0.680$ at
        $W/L=1.67$, $1.262$ at $W/L=5$; domain width must be matched to
        experiment.
  \item Scaling law $t_{\mathrm{crit}} \propto \sqrt{\Gc/\ell}$ provides a
        direct experimental handle: measuring fracture time under controlled
        growth rates constrains the ratio $\Gc/\ell$.
\end{itemize}

\noindent\textbf{Open points:}
\begin{enumerate}
  \item Spatial heterogeneity in $\alpha(x,z,t)$ to nucleate multiple
        substrate cracks and predict crack spacing.
  \item Tensile/compressive split of $\psiE$ (Miehe et al.\ 2010) for mixed-mode
        fracture near the substrate corners.
  \item Experimental calibration of $\Gc$ and $\ell$ from biofilm
        detachment assays (AFM-based adhesion energy).
\end{enumerate}

% ─────────────────────────────────────────────────────────────────────────────
\section*{Code Availability}

All simulation and plotting scripts are in the
\texttt{pde-fem-biofilm} repository (\href{https://github.com/keisuke58/pde-fem-biofilm}{keisuke58/pde-fem-biofilm}):
\begin{itemize}
  \item \texttt{phase\_field\_at2\_2d\_sparse.py} --- solver (prescribed + coupled modes)
  \item \texttt{plot\_at2\_thesis.py} --- produces \cref{fig:coupled}
  \item \texttt{plot\_param\_study\_thesis.py} --- produces \cref{fig:params}
\end{itemize}
Run environment: Python~3.12, \texttt{scipy}~1.13, \texttt{matplotlib}~3.9,
TeX~Live~2025 (lmodern).

\end{document}
